\documentclass[10pt]{article}

% Specify the margins and text size.
\setlength{\textwidth}{6.6in}
\setlength{\textheight}{9.8in}
\setlength{\oddsidemargin}{0pt}
\setlength{\evensidemargin}{0pt}
\setlength{\topmargin}{0pt}
\setlength{\hoffset}{.05in}
\setlength{\voffset}{-1in}

\setlength{\parskip}{5pt}
\setlength{\parindent}{0pt}

% Load some fonts and symbol packages
\usepackage{latexsym}
\usepackage{pifont}       % contains 'star' symbol for counterinsurgency handbook title
\usepackage{yfonts} 
\usepackage{amsmath}
\usepackage{amsfonts}

\usepackage{graphicx}     % actually, this is loaded by pstricks
\usepackage[T1]{fontenc}
\usepackage{ifthen}
\usepackage{pstricks,pst-grad,pst-func,pst-text,pst-node,multido,pst-plot,calc,pst-3dplot}
%\usepackage[all]{xy}
%\usepackage{animate}

% The hyperref package inserts links into the pdf.
\definecolor{MyLinkColor}{rgb}{.1,.2,1}
\definecolor{MyCiteColor}{rgb}{.1,1,.2}
\definecolor{MyURLColor}{rgb}{.4,.4,.4}
\usepackage[backref=true,pagebackref=false,hyperindex,colorlinks=true,
  linkcolor=MyLinkColor,urlcolor=MyURLColor]{hyperref}


% The tweaklist package is something I found on the web.  It provides a simple interface
% for making changes to spacing used in the itemize and enumerate environments.  Comment
% this out if you don't care to use tweaklist.
\usepackage{tweaklist}
\renewcommand{\itemhook}{\setlength{\parskip}{2pt}\setlength{\parsep}%
{1pt}\setlength{\topsep}{0pt}\setlength{\itemsep}{0pt}}

\newcommand{\U}{\underline{\hspace{5pt}}}

\usepackage{listings}
\newcounter{EX}\setcounter{EX}{1}
\newcommand{\EXERCISE}{\arabic{EX}.\stepcounter{EX} }

\begin{document}
\pagestyle{empty}
\lstset{language=R, showspaces=false, showstringspaces=false}

\href{http://www.su.edu}{\includegraphics[height=1.75cm]{sulogo.eps}}
\vspace{-1.69cm}

{{\ }\hfill\small
\begin{tabular}{cl}
& Math 207\\
& Introduction to Statistics\\
\end{tabular}
}
\setlength{\baselineskip}{1.05\baselineskip}

\begin{center}
\textbf{\large  Exercises:  The Normal Curve}
\end{center}
\newcommand{\Z}{\hphantom{0}}

\begin{center}
\begin{tabular}{ccc}
\begin{pspicture}(-3,-2.5)(4,1.75)
\psframe[linewidth=0.02](-3.8,-1.1)(4.2,1.8)
\psset{yunit=3,linewidth=0.02,arrowsize=3pt 2}
%
   \pscustom[fillstyle=solid,fillcolor=lightgray,sigma=1,linestyle=none]{
     \psGauss{-3.5}{1.25}
     \psline(1.25, 3.490343)(1.25,0)(-3.5,0) }
   \psGauss[sigma=1,linecolor=blue,linewidth=0.5pt]{-3.5}{3.5}
\psline(1.25,0.1826)(1.25,-0.22)   \rput[t](1.25,-0.25){\scriptsize 1.25}
\rput[l](1.5,0.3){\scriptsize\texttt{pnorm}$(1.25)\approx 0.894$}
\psline{->}(1.45,0.29)(0.5,0.2)
%
\psline(-3.5,0)(3.5,0)  
  \psline(-3,0)(-3,-0.04) \rput[t](-3,-0.07){\scriptsize -3\hphantom{-}}
  \psline(-2,0)(-2,-0.04) \rput[t](-2,-0.07){\scriptsize -2\hphantom{-}}
  \psline(-1,0)(-1,-0.04) \rput[t](-1,-0.07){\scriptsize -1\hphantom{-}}
  \psline(0,0)(0,-0.04)   \rput[t](0,-0.07){\scriptsize 0}
  \psline(1,0)(1,-0.04)   \rput[t](1,-0.07){\scriptsize 1}
  \psline(2,0)(2,-0.04)   \rput[t](2,-0.07){\scriptsize 2}
  \psline(3,0)(3,-0.04)   \rput[t](3,-0.07){\scriptsize 3}
  \rput[l](3.6,0){\scriptsize $x$}
\psline(0,0)(0,0.5)
  \psline(-0.12,0.5)(0,0.5)    \rput[r](-0.21,0.5){\scriptsize $0.5$}
  \psline(-0.12,0.25)(0,0.25)  \rput[r](-0.21,0.25){\scriptsize $0.25$}
\end{pspicture}
&
&
\begin{pspicture}(-4,-2.5)(3,1.75)
\psframe[linewidth=0.02](-3.8,-1.1)(4.2,1.8)
\psset{yunit=3,linewidth=0.02,arrowsize=3pt 2}
%
   \pscustom[fillstyle=solid,fillcolor=lightgray,sigma=1,linestyle=none]{
     \psGauss{-0.5}{1.25}
     \psline(1.25, 3.490343)(1.25,0)(-0.5,0)(-0.5,0.352065)}
   \psGauss[sigma=1,linecolor=blue,linewidth=0.5pt]{-3.5}{3.5}
\psline(1.25,0.1826)(1.25,-0.22)   \rput[t](1.25,-0.25){\scriptsize 1.25}
\psline(-0.5,0.352065)(-0.5,-0.22)   \rput[t](-0.5,-0.25){\scriptsize $-0.5$}
\rput[l](0.8,0.45){\scriptsize$\mbox{\texttt{pnorm}}(1.25)-\mbox{\texttt{pnorm}}(-0.5)$}
\rput[l](2.5,0.33){\scriptsize$\approx 0.586$}
\psline{->}(2.1,0.32)(0.5,0.1)
%
\psline(-3.5,0)(3.5,0)  
  \psline(-3,0)(-3,-0.04) \rput[t](-3,-0.07){\scriptsize -3\hphantom{-}}
  \psline(-2,0)(-2,-0.04) \rput[t](-2,-0.07){\scriptsize -2\hphantom{-}}
  \psline(-1,0)(-1,-0.04) \rput[t](-1,-0.07){\scriptsize -1\hphantom{-}}
  \psline(0,0)(0,-0.04)   \rput[t](0,-0.07){\scriptsize 0}
  \psline(1,0)(1,-0.04)   \rput[t](1,-0.07){\scriptsize 1}
  \psline(2,0)(2,-0.04)   \rput[t](2,-0.07){\scriptsize 2}
  \psline(3,0)(3,-0.04)   \rput[t](3,-0.07){\scriptsize 3}
  \rput[l](3.6,0){\scriptsize $x$}
\psline(0,0)(0,0.5)
  \psline(-0.12,0.5)(0,0.5)    \rput[r](-0.21,0.5){\scriptsize $0.5$}
  \psline(-0.12,0.25)(0,0.25)  \rput[r](-0.21,0.25){\scriptsize $0.25$}
\end{pspicture}\\[-30pt]
%%
%%
\begin{pspicture}(-3,-2.5)(4,1.75)
\psframe[linewidth=0.02](-3.8,-1.1)(4.2,1.8)
\psset{yunit=3,linewidth=0.02,arrowsize=3pt 2}
%
   \pscustom[fillstyle=solid,fillcolor=lightgray,sigma=1,linestyle=none]{
     \psGauss{3.5}{1.25}
     \psline(1.25, 3.490343)(1.25,0)(3.5,0) }
   \psGauss[sigma=1,linecolor=blue,linewidth=0.5pt]{-3.5}{3.5}
\psline(1.25,0.1826)(1.25,-0.22)   \rput[t](1.25,-0.25){\scriptsize 1.25}
\rput[l](0.8,0.38){\scriptsize$1-\mbox{\texttt{pnorm}}(1.25)\approx 0.106$}
\psline{->}(1.45,0.29)(1.5,0.04)
%
\psline(-3.5,0)(3.5,0)  
  \psline(-3,0)(-3,-0.04) \rput[t](-3,-0.07){\scriptsize -3\hphantom{-}}
  \psline(-2,0)(-2,-0.04) \rput[t](-2,-0.07){\scriptsize -2\hphantom{-}}
  \psline(-1,0)(-1,-0.04) \rput[t](-1,-0.07){\scriptsize -1\hphantom{-}}
  \psline(0,0)(0,-0.04)   \rput[t](0,-0.07){\scriptsize 0}
  \psline(1,0)(1,-0.04)   \rput[t](1,-0.07){\scriptsize 1}
  \psline(2,0)(2,-0.04)   \rput[t](2,-0.07){\scriptsize 2}
  \psline(3,0)(3,-0.04)   \rput[t](3,-0.07){\scriptsize 3}
  \rput[l](3.6,0){\scriptsize $x$}
\psline(0,0)(0,0.5)
  \psline(-0.12,0.5)(0,0.5)    \rput[r](-0.21,0.5){\scriptsize $0.5$}
  \psline(-0.12,0.25)(0,0.25)  \rput[r](-0.21,0.25){\scriptsize $0.25$}
\end{pspicture}
&
&
\begin{pspicture}(-4,-2.5)(3,1.75)
\psframe[linewidth=0.02](-3.8,-1.1)(4.2,1.8)
\psset{yunit=3,linewidth=0.02,,arrowsize=3pt 2}
%
   \pscustom[fillstyle=solid,fillcolor=lightgray,sigma=1,linestyle=none]{
     \psGauss{-3.5}{0.6744898}
     \psline(0.6744898, 0.3177766)(0.6744898,0)(-3.5,0) }
   \psGauss[sigma=1,linecolor=blue,linewidth=0.5pt]{-3.5}{3.5}
\psline(0.6744898,0.3177766)(0.6744898,-0.22)   
   \rput[tl](0.61,-0.23){\scriptsize\texttt{qnorm}$(0.75)\approx 0.67$}
\rput[l](1.5,0.3){\scriptsize$\mbox{area}=0.75$}
\psline{->}(1.45,0.29)(0.25,0.15)
%
\psline(-3.5,0)(3.5,0)  
  \psline(-3,0)(-3,-0.04) \rput[t](-3,-0.07){\scriptsize -3\hphantom{-}}
  \psline(-2,0)(-2,-0.04) \rput[t](-2,-0.07){\scriptsize -2\hphantom{-}}
  \psline(-1,0)(-1,-0.04) \rput[t](-1,-0.07){\scriptsize -1\hphantom{-}}
  \psline(0,0)(0,-0.04)   \rput[t](0,-0.07){\scriptsize 0}
  \psline(1,0)(1,-0.04)   \rput[t](1,-0.07){\scriptsize 1}
  \psline(2,0)(2,-0.04)   \rput[t](2,-0.07){\scriptsize 2}
  \psline(3,0)(3,-0.04)   \rput[t](3,-0.07){\scriptsize 3}
  \rput[l](3.6,0){\scriptsize $x$}
\psline(0,0)(0,0.5)
  \psline(-0.12,0.5)(0,0.5)    \rput[r](-0.21,0.5){\scriptsize $0.5$}
  \psline(-0.12,0.25)(0,0.25)  \rput[r](-0.21,0.25){\scriptsize $0.25$}
\end{pspicture}
\end{tabular}\vspace{-.6in}
\end{center}


1. Sketch then, without using a calculator, estimate the specified area under the normal curve.

\begin{center}
\begin{tabular}{lcl}
a) $z>2$ & \hspace{3in} & b) $z>2$\\[.6in]
c) $-1<z<2$ &           & d) $1<z<2$\\[.6in]
e) $z<10$   &           & f) $z<-7$\hphantom{\hspace{1in}}\\[.6in]
g) $-3<z<1$ &           & h) $|z|<1$\\[.6in]
\end{tabular}
\end{center}

2. Use the \texttt{pnorm} function in \texttt{R} to find the specified areas under the normal curve.
Write down the commands that you use.
\begin{center}
\begin{tabular}{lcl}
a) $z<-1.5$ & \hspace{3in} & b) $z>1.8$\\[.6in]
c) $z<7.75$   &           & d) $z<-5.3$\hphantom{\hspace{1in}}\\[.6in]
e) $-1.25<z<2.25$ &           & f) $1.8<z<2.0$\\
\end{tabular}
\end{center}
\vfill
\eject

3. Suppose that math SAT scores average out to 520 with a standard deviation of 20.
Assume that the histogram of scores follows a normal distribution.

\hspace{10pt} a) About what percent of scores were between 500 and 540?
\vspace{.8in}

\hspace{10pt} b) About what percent of scores were below 560?
\vspace{.8in}

\hspace{10pt} c) About what percent of scores were over 580?
\vspace{.8in}

\hspace{10pt} d) What score is better than 90\% of all scores? (Hint:  \texttt{qnorm}.)
\vspace{1in}

4. Errors in the measurement of a certain check weight average out to $0.0065$ kg plus or minus
about $0.0002$~kg.  Assume that the histogram of measurement errors follows a normal 
distribution.

\hspace{10pt} a) What is the chance of getting an error measurement greater than 
$0.0068$?
\vspace{.8in}

\hspace{10pt} b) What is the chance of getting an error measurement greater than 
between $0.0064$ and $0.0066$?
\vspace{.8in}

\hspace{10pt} c) About what percent of measurements  were over $0.0070$?
\vspace{.8in}

\hspace{10pt} d) What measurement is larger than 90\% of all measurements? (Hint:  \texttt{qnorm}.)
\vfill
\eject

\end{document}



